\section{Análisis de anonimización}

Los datos tabulares no contienen datos personales: No hay nombres de trabajadores, operadores ni clientes. Solo se identifican vacas y grupos de manejo. No obstante, es información interna del CAETEC y no deben hacerse públicos.

Donde sí encontramos datos personales es en los metadatos de los archivos y en las fotos. En Google Drive, cada archivo muestra a su dueño con su matrícula o correo, y las fotos pueden guardar la ubicación GPS y el modelo del teléfono con el que se tomaron. Además, en algunas fotos podrían aparecer trabajadores o estudiantes. En la siguiente tabla se describe cada atributo, si identifica a una persona y qué se hará con él.

\begin{table}[htbp]
  \centering
  \small
  \begin{tabularx}{\textwidth}{>{\raggedright\arraybackslash}X>{\raggedright\arraybackslash}p{0.14\textwidth}>{\raggedright\arraybackslash}X>{\raggedright\arraybackslash}X}
    \toprule
    Atributo & ¿Identifica a una persona? & Riesgo & Acción \\
    \midrule
    ID de vaca & No & Información productiva y reproductiva del CAETEC & Se conserva, pero no se publica \\
    Grupo y robot (“Grupo 6”, “VMS 3”) & No & Bajo & Se conserva \\
    Fechas y horas de ordeño & No & Bajo & Se conserva \\
    Dueño de los archivos en Drive (matrícula, correo) & Sí & Datos personales de estudiantes y profesores & No se copia a los datos ni al repositorio \\
    Metadatos de las fotos (GPS, dispositivo) & Posible & Ubicación del establo y teléfono de un estudiante & Se eliminan antes de subir las fotos \\
    Personas en las fotos & Posible & Imagen de trabajadores o estudiantes & Se recorta o se descarta la foto \\
    \bottomrule
  \end{tabularx}
\end{table}

Todavía falta revisar el contenido del archivo “Datos 2026 TEC-UAQ-UNAM”, ya que Drive no permite previsualizarlo sin descargarlo.
